\chapter{Pareja de Planck, tensi\'on universal y compliancia gravitatoria}
\label{chap:pareja-planck-respuesta}

\section{Pareja dimensional de Planck \((Q,\ell)\): \(Q\ell=\hbar c\), \(\ell/Q=G/c^4\) y reconstrucción unívoca}

Sea \(t_0\) la unidad cronol\'ogica de la secci\'on CT108 y sea \(c\) la
secci\'on causal seleccionada por el lector constitutivo. Definimos

\[
t_P=\frac{54}{\pi}t_0,
\qquad
\ell=\ell_P=ct_P=\frac{54}{\pi}ct_0,
\]

y el cuanto energ\'etico

\[
Q=E_P=\frac{\hbar}{t_P}=\frac{h}{108t_0}.
\]

La periodicidad \(T_{108}=108t_0\) verifica

\[
T_{108}=2\pi t_P,
\qquad
QT_{108}=h.
\]

De estas definiciones se deduce la primera identidad de la pareja:

\begin{equation}
\boxed{Q\ell=\hbar c.}
\label{eq:rm-planck-product}
\end{equation}

La realizaci\'on gravitatoria circular determina

\[
G=\left(\frac{54}{\pi}\right)^2\frac{c^5t_0^2}{\hbar}.
\]

Al sustituir \(\ell=(54/\pi)ct_0\) y \(Q=\hbar/t_P\), se obtiene la
segunda identidad:

\begin{equation}
\boxed{\frac{\ell}{Q}=\frac{G}{c^4}.}
\label{eq:rm-planck-quotient}
\end{equation}

Las ecuaciones \eqref{eq:rm-planck-product} y
\eqref{eq:rm-planck-quotient} son algebraicamente independientes como
relaciones de producto y cociente. Reconstruyen de forma un\'ivoca

\[
Q=\sqrt{\frac{\hbar c^5}{G}},
\qquad
\ell=\sqrt{\frac{G\hbar}{c^3}}.
\]

\begin{theorem}[Clausura de la pareja de Planck]
\label{thm:rm-planck-pair}
En la realizaci\'on circular CT108, la pareja positiva \((Q,\ell)\) es la
\'unica soluci\'on de
\[
Q\ell=\hbar c,
\qquad
\ell/Q=G/c^4.
\]
Satisface, adem\'as,
\[
GQ=c^4\ell,
\qquad
GQ^2=\hbar c^5,
\qquad
G\hbar=c^3\ell^2.
\]
\end{theorem}

\begin{proof}
La positividad permite multiplicar y dividir las dos ecuaciones sin
ambig\"uedad de signo. La primera igualdad adicional procede de
\(\ell/Q=G/c^4\); la segunda resulta al multiplicarla por \(Q^2\), y la
tercera al multiplicar \(\ell/Q=G/c^4\) por
\(Q\ell=\hbar c\). La reconstrucci\'on por ra\'ices positivas demuestra la
unicidad.
\end{proof}

\section{Reciprocidad entre tensión y compliancia de Planck y conservación del producto de acción}

La pareja define dos respuestas rec\'iprocas:

\begin{equation}
\mathscr T_P:=\frac{Q}{\ell}=\frac{c^4}{G},
\qquad
\mathscr C_P:=\frac{\ell}{Q}=\frac{G}{c^4}.
\label{eq:rm-tension-compliance}
\end{equation}

La primera posee dimensi\'on de fuerza; la segunda, dimensi\'on de
compliancia gravitatoria. Satisfacen

\[
\mathscr T_P\mathscr C_P=1.
\]

La ecuaci\'on de Einstein se escribe entonces como una ley constitutiva:

\begin{equation}
\boxed{
G_{\mu\nu}+\Lambda g_{\mu\nu}
=8\pi\mathscr C_P T_{\mu\nu}.}
\label{eq:rm-einstein-compliance}
\end{equation}

La materia aporta tensi\'on energ\'etica y la geometr\'ia responde mediante
\(\mathscr C_P\). La constante \(G\) publica la capacidad de curvatura por
unidad de tensi\'on en la carta causal elegida.

En el r\'egimen newtoniano, para una densidad energ\'etica
\(\varepsilon\),

\[
\nabla^2\!\left(\frac{\Phi}{c^2}\right)
=4\pi\mathscr C_P\varepsilon.
\]

La acci\'on de Einstein--Hilbert queda normalizada por el \`area estructural:

\begin{equation}
\frac{S_{\rm EH}}{\hbar}
=\frac{1}{16\pi\ell^2}
\int_{\mathcal M}(R-2\Lambda)\sqrt{-g}\,\mathrm d^4x.
\label{eq:rm-eh-area}
\end{equation}

La respuesta escalar y la respuesta tensorial son publicaciones diferentes
de la misma compliancia. La primera utiliza \(4\pi\mathscr C_P\); la segunda,
\(8\pi\mathscr C_P\). Al insertar CT108,

\[
2\pi\mathscr C_P=108\frac{ct_0}{Q},
\qquad
4\pi\mathscr C_P=216\frac{ct_0}{Q},
\qquad
8\pi\mathscr C_P=432\frac{ct_0}{Q}.
\]

La progresi\'on \(108\to216\to432\) se deduce de los factores geom\'etricos.
La coincidencia del \(432\) con la multiplicidad de la regi\'on transversal
de \(\pi\) constituye una compatibilidad estructurada; su promoci\'on requiere
un mapa de incidencia que transporte la multiplicidad geneal\'ogica al tensor
de Einstein.

\section{Cuanto informacional ternario \(\log 3\) y capacidad de la celda qutrit}

La relaci\'on termodin\'amica de la secci\'on CT108 es

\begin{equation}
\boxed{Q=k_B\Theta_{\rm clk}\log3.}
\label{eq:rm-planck-trit}
\end{equation}

En esta ecuaci\'on, \(\log3\) es la entrop\'ia de una alternativa ternaria
equiprobable; \(\Theta_{\rm clk}\) fija la carta t\'ermica y \(Q\) publica la
energ\'ia del ciclo. La combinaci\'on con el teorema
\ref{thm:rm-planck-pair} produce

\begin{equation}
\boxed{
Gk_B\Theta_{\rm clk}\log3=c^4\ell.}
\label{eq:rm-gravity-information}
\end{equation}

La gravedad transforma el cuanto informacional del reloj en una longitud de
respuesta geom\'etrica. Equivalentemente,

\[
\ell=\mathscr C_P k_B\Theta_{\rm clk}\log3.
\]

Esta igualdad separa el contenido adimensional \(\log3\), el marco
termom\'etrico \(k_B\Theta_{\rm clk}\) y la compliancia geom\'etrica.

\section{Covariancia de los lectores bidireccionales bajo inversión de hoja y conjugación temporal}

Sean \(\hbar_{\rm pre}\) y \(\hbar_{\rm ret}\) las dos secciones orientadas
de la recta de acci\'on, relacionadas por un cociente positivo
\(R_{\rm act}\):

\[
\frac{\hbar_{\rm pre}}{\hbar_{\rm ret}}=R_{\rm act}.
\]

La invariancia de \(\ell\) obliga a

\[
\frac{G_{\rm pre}}{G_{\rm ret}}=R_{\rm act}^{-1}.
\]

Por tanto,

\begin{equation}
G_{\sigma}\hbar_{\sigma}=c^3\ell^2
\qquad(\sigma\in\{\mathrm{pre},\mathrm{ret}\}).
\label{eq:rm-two-way-area}
\end{equation}

La acci\'on conserva la orientaci\'on de la ruta; el producto con la respuesta
gravitatoria conserva el \`area. Esta covariancia convierte a \(\ell^2\) en
el invariante geom\'etrico de la inversi\'on bidireccional.

\section{Vac\'io constitutivo y cierre causal}

El operador de vac\'io proporciona

\[
c_{\rm vac}^{-2}=\mu\varepsilon.
\]

La pareja de Planck proporciona una segunda secci\'on causal:

\[
c_{\rm int}=\frac{Q\ell}{\hbar}.
\]

Definimos el escalar de costura

\begin{equation}
\Xi_c=\frac{Q\ell}{\hbar}\sqrt{\mu\varepsilon}
=\frac{c_{\rm int}}{c_{\rm vac}}.
\label{eq:rm-causal-seam}
\end{equation}

La igualdad \(\Xi_c=1\) es necesaria y suficiente para que ambos lectores
publiquen el mismo cono causal. Bajo esta costura,

\begin{equation}
\boxed{
G=\left(\frac{54}{\pi}\right)^2
\frac{t_0^2}{\hbar(\mu\varepsilon)^{5/2}}.}
\label{eq:rm-g-vacuum-product}
\end{equation}

\begin{proof}
La relaci\'on \(c=(\mu\varepsilon)^{-1/2}\) da
\(c^5=(\mu\varepsilon)^{-5/2}\). Su sustituci\'on en el selector circular de
\(G\) produce \eqref{eq:rm-g-vacuum-product}.
\end{proof}

El producto \(\mu\varepsilon\) gobierna el cono causal y la respuesta
gravitatoria; el cociente \(\mu/\varepsilon=Z^2\) gobierna impedancia y
orientaci\'on. Ambos contenidos proceden del mismo operador bicapa y conservan
funciones distintas.

\section{Factorización machiana de la respuesta inercial local mediante la traza del cierre global: \(\mathcal I_v=\overline{\mathcal I}_v\circ b\)}

Una energ\'ia local \(E\) determina la coordenada adimensional

\[
\xi_E(E)=\frac{E}{Q}.
\]

El acoplamiento gravitatorio de dos energ\'ias es

\begin{equation}
\alpha_G(E_1,E_2)=\frac{E_1E_2}{Q^2}.
\label{eq:rm-mach-coupling}
\end{equation}

La referencia \(Q\) pertenece al reloj global de cierre. La inercia local se
eval\'ua, por tanto, frente a una escala producida por la genealog\'ia completa.
Esta dependencia no postula una acci\'on instant\'anea del universo sobre el
cuerpo: establece una factorizaci\'on operatoria de la respuesta local por la
clase global de cierre.

\begin{proofstatusbox}
Las identidades de la pareja, la tensi\'on, la compliancia y la covariancia
\(G\hbar\) son exactas en la realizaci\'on circular CT108 declarada. La
identificaci\'on \(c_{\rm vac}=c_{\rm int}\) es una interfaz f\'isica tipada;
su contenido se concentra en el mapa de rectas causales cuya evaluaci\'on es
\(\Xi_c=1\).
\end{proofstatusbox}

\begin{falsifierbox}
Refutan la clausura de este cap\'itulo: una pareja positiva distinta que
satisfaga simult\'aneamente las dos ecuaciones rectoras; la p\'erdida de
\(G\hbar=c^3\ell^2\) bajo la inversi\'on de acci\'on; una secci\'on de vac\'io
que no cumpla \(c^{-2}=\mu\varepsilon\); o un transporte de marco para el que
\(\Xi_c\ne1\).
\end{falsifierbox}
